\documentclass[../../../../main.tex]{subfiles}
\begin{document}

\begin{flushright}
\footnotesize\textit{【自動文字起こし・要確認】}
\end{flushright}

[解] (1) 1回目で勝負がキマルとき. $1 \le r \le s+1$ であるから以下このもとで考える. $\cdots \text{\textcircled{1}}$
(1) 1回目に勝者が決まる確率は, $s=0$の時 $p_1 = 1$, $s \ge 1$の時, $r=1$なら $p_1 = \frac{r}{r+s}$, $r \ge 2$なら
$$ p_1 = \frac{s}{r+s} \cdot \frac{s-1}{r+s-1} \cdots \frac{s-(r-2)}{r+s-(r-2)} \cdot \frac{r}{r+s-(r-1)} $$
だから $r=1,2,3,\dots,s$の時, 一部を$A$として,
\begin{align*}
p_r - p_{r+1} &= A \cdot \frac{r}{r+s-(r-1)} - A \cdot \frac{s-(r-1)}{r+s-(r-1)} \cdot \frac{r}{r+s-r} \\
&= A r \frac{r-1}{(r+s-r)(r+s-r+1)} \ge 0 \quad \cdots *1
\end{align*}
$r=1$の時
$$ p_1 - p_2 = \frac{r}{r+s} - \frac{s}{r+s} \cdot \frac{r}{r+s-1} \ge 0 \ (\because s \ge 1) \cdots *2 $$
以上から, $p_r \ge p_{r+1}$ である. (終)

以上から $p_A \ge \frac{1}{2}$ であり, 等号が成立するのは
$1^\circ$の時の, $p_k \ge p_{k+1} \ (k=1,2 \dots 2s-1)$ とした式全てで
等号が成立する時である. この時まず$*1$での成立が必要で,
$*2$の等号が成立し, $r \ge 1, s \ge 0$から $r=1$が必要. この時 $s=1$なら坊.
$s \ge 2$の時, $*2$の等号も$r$によらず成立し, $r=1$である.
よって, もとめる条件は $r=1, s \in \text{odd}$ \hfill \#

(2) Aが勝者となるのは, $2k-1 \ (k \in \mathbb{N}, \ 1 \le 2k-1 \le s+1)$ 回目に
はじめて赤玉を引く確率で, これを$P_A$, Bが勝者となる確率を$P_B$とおくと
$1^\circ \ s=2s'-1 \ (s' \in \mathbb{N})$の時
$$
\begin{cases}
P_A = p_1 + p_3 + \dots + p_{2s'-1} \\
P_B = p_2 + p_4 + \dots + p_{2s'}
\end{cases}
$$
であり, (1)から $p_1 \ge p_2, \dots, p_{2s'-1} \ge p_{2s'}$ とした式を辺々たしあわせて
$$ P_A \ge P_B $$
これと$P_A + P_B = 1$から, $P_A \ge \frac{1}{2}$

$2^\circ \ s=2s' \ (s \in \mathbb{N}_0)$の時
$$
\begin{cases}
P_A = p_1 + p_3 + \dots + p_{2s'-1} + p_{2s'+1} \\
P_B = p_2 + p_4 + \dots + p_{2s'}
\end{cases}
$$
であり, 同様に $p_1 \ge p_2, \dots$ などとした式と $p_{2s'+1} \ge 0$から,
$$ P_A \ge P_B \quad \therefore P_A \ge \frac{1}{2} \ (\because P_A + P_B = 1) $$

$3^\circ \ s=0$の時
Aが1回目で勝利するから $P_A = 1$

\end{document}
